\documentclass[a4paper, twoside]{report}

%% Language and font encodings
\usepackage[english]{babel}
\usepackage[utf8x]{inputenc}
\usepackage[T1]{fontenc}
\usepackage{mathrsfs}

%% Sets page size and margins
\usepackage[a4paper,top=3cm,bottom=2cm,left=3cm,right=3cm,marginparwidth=1.75cm]{geometry}

\usepackage{amsthm}
\usepackage{graphicx}
\usepackage[]{listings}
\usepackage{booktabs}
\usepackage[colorinlistoftodos]{todonotes}
\usepackage[colorlinks=true, allcolors=blue]{hyperref}
\usepackage{comment}
\usepackage[numbers]{natbib}
\usepackage{fancyref}
\usepackage{wrapfig}

%% Math
\usepackage{mathtools}
\usepackage{stmaryrd}

\usepackage{title/titlestructure}

%% Utils
\usepackage{background/utils}


\title{A Deeper Dive into Relational Algebra by Way of Adjunctions}
\author{Matteo Bongiovanni}
\date{January 2023}
\supervisor{Dr.\ Nicolas Wu}
\secondmarker{Dr.\ Steffen van Bakel}
\documenttype{MEng Individual Project}
\university{Imperial College London}
\department{Department of Computing}
\degree{Joint Mathematics and Computing}

\begin{document}
\begin{comment}
General Information
The project report (occasionally referred to as a thesis) is an extremely important aspect of the project. It serves to show what you have achieved and should demonstrate that:

You understand the wider context of computing by relating your choice of project, and the approach you take, to existing products or research.
You can apply the theoretical and practical techniques taught in the course to the problem you are addressing and that you understand their relevance to the wider world of computing.
You are capable of objectively criticising your own work and making constructive suggestions for improvements or further work based on your experiences so far.
As a computing professional, you operate ethically and can explain your thinking and working processes clearly and concisely to third parties who may not be experts in the field in which you are working.
Most of the project assessors will not have followed the project throughout and will only have a short time to listen to a presentation or see a demonstration. For this reason they will rely heavily on the report to judge the project.

Many students underestimate the importance of the report and make the mistake of thinking that top marks can be achieved simply for producing a good product. This is fundamentally not the case and many projects have been graded well below their potential because of an indifferent or poor write-up. In order to get the balance right you should consider that the aim of the project is to produce a good report and that software, hardware, theory etc. that you developed during the project are merely a means to this end. Don't make the mistake of leaving the write-up to the last minute. Ideally you should produce the bulk of the report as you go along and use the last week or two to bring it together into a coherent document.

It is helpful to get feedback from your supervisor about your project report, but supervisors cannot be expected to look at documents written at the last minute or at more than one chapter at a time. Allow plenty of time for this.

Project Report Length
The nominal maximum report length is 60 pages, plus appendices. Shorter reports may be appropriate for many projects and longer reports can be submitted with prior approval (see below).

The page count includes:

title page
abstract
dedications
table of contents
The page count does not include:

bibliography
appendices (but the report should make sense without the appendices)
When writing up your work be reminded of the fact that a 60-page report isn't necessarily better than a 30-page report, in the same way that a 10,000 line implementation isn't necessarily better 5,000 line one. Conciseness, clarity and elegance are invaluable qualities in report writing, just as they are in programming. Indeed, you will lose credit if it is evident to the assessors that the report has been padded with superfluous content in order to make it appear more substantial than the work justifies. The assessors greatly value quality over quantity.

It is important to understand that the report is assessed on the quality of the technical writing. Here this means your ability to document your ideas, design, implementation, evaluation, conclusions etc. succinctly, using appropriate scientific language, as might a researcher when writing up their findings in an technical magazine, academic journal or conference proceedings.

If you have good reason to use more than 60 pages you should consult your supervisor and second marker and obtain their approval in advance - they will be required to read and assess the whole report, whereas other assessment team members may only scan key sections.

Project Layout
The physical layout and formatting of the report is also important. A tidy, well laid out and consistently formatted document makes for easier reading and is suggestive of a careful and professional attitude towards its preparation. Both LaTeX and MS Word will allow you to produce a cleanly formatted document. Many supervisors will advise you to use LaTeX as this solves most of the formatting problems for you. If your report is to contain a substantial number of mathematical formulae you are strongly advised to use LaTeX.

Project Report Structure
The exact structure of your project (the chapter titles etc) is up to you. Discuss with your supervisor what structure works best for your project.

A project report will usually contain a number of chapters (Chapter 1, Chapter 2 etc). Each chapter contains Sections (2.1, 2.2 etc) and may also contain subsections (2.1.1, 2.1.2) etc. Try to avoid too many levels of subheading - three is usually sufficient.

Projects will usually contain all of the following elements, although these may not exactly map to chapters.
\end{comment}

% Front matter
\begin{comment}
This should include the project title and the name of the author of the report. You can also list the name of your supervisor if you wish. 
\end{comment}

\begin{titlepage}

\newcommand{\HRule}{\rule{\linewidth}{0.5mm}} % Defines a new command for the horizontal lines, change thickness here

%----------------------------------------------------------------------------------------
%	LOGO SECTION
%----------------------------------------------------------------------------------------

\includegraphics[width=8cm, draft=false]{report/title/logo.eps}\\[1cm] % Include a department/university logo - this will require the graphicx package
 
%----------------------------------------------------------------------------------------

\center\ % Center everything on the page

%----------------------------------------------------------------------------------------
%	HEADING SECTIONS
%----------------------------------------------------------------------------------------

\makeatletter
\textsc{\LARGE \@documenttype}\\[1.5cm]
\textsc{\Large \@university}\\[0.5cm]
\ifdefined\@degree\
\textsc{\large \@degree}\\[0.5cm]
\fi
\textsc{\large \@department}\\[0.5cm]
\makeatother

%----------------------------------------------------------------------------------------
%	TITLE SECTION
%----------------------------------------------------------------------------------------
\makeatletter
\HRule\\[0.4cm]
{ \huge \bfseries \@title}\\[0.4cm] % Title of your document
\HRule\\[1.5cm]
 
%----------------------------------------------------------------------------------------
%	AUTHOR SECTION
%----------------------------------------------------------------------------------------

\begin{minipage}{0.4\textwidth}
\begin{flushleft} \large
\emph{Author:}\\
\@author\ % Your name
\end{flushleft}
\end{minipage}
~\begin{minipage}{0.4\textwidth}
\begin{flushright} \large
\emph{Supervisor:} \\
\@supervisor\\[1.2em]

\ifdefined\@secondmarker\
\emph{Second Marker:} \\
\@secondmarker\
\fi

\end{flushright}
\end{minipage}\\[2cm]

%----------------------------------------------------------------------------------------
%	DATE SECTION
%----------------------------------------------------------------------------------------

{\large \@date}\\[2cm] % Date, change the \today to a set date if you want to be precise

\vfill % Fill the rest of the page with whitespace
\makeatother

\end{titlepage}
\begin{comment}
The abstract is a very brief summary of the report's contents. It should be no more than half a page long. Somebody unfamiliar with your project should have a good idea of what it's about having read the abstract alone and will know whether it will be of interest to them. Note that the abstract is a summary of the entire project including its conclusions. A common mistake is to provide only introductory elements in the abstract without saying what has been achieved.
\end{comment}


\begin{abstract}
In order to evaluate alternative equijoin implementation in Haskell, this paper
introduced the \relation{JOINBENCH} relation and surrounding methodology and
tooling. The \relation{JOINBENCH} relation is a scheme designed for the
synthesis of data optimised to help evaluate the performance of equijoins in a
variety of scenarios. Furthermore, this paper presents a low-level library that
helps users customise and define their own synthetic data sources for future
benchmarking purposes.

In their distinguished paper
\relalg{}~\cite{RelationalAlgebraByWayOfAdjunctions} it was noted that the
monadic structure of bulk types can help explain most of relational algebra.
Using this structure, the authors designed a new method to facilitate
the use of monad comprehensions in an efficient implementation of equijoins of
relational databases. This project presents an implementation of such a system
and an evaluation of the performance such query optimisations carry using the
bespoke tooling described above. 
\end{abstract}

\begin{comment}
Table of contents:
This should list the main chapters, sections and subsections of your report. Choose self-explanatory chapter and section titles and use double spacing for clarity. If possible you should include page numbers indicating where each chapter/section begins. 
\end{comment}
\tableofcontents

% Main matter 
\chapter{Introduction}
\begin{comment}
The introduction should summarise the subject area, the specific problem you are addressing, including key ideas for their solution, together with a summary of the project's main contributions. When detailing the contributions it is helpful to provide forward references to the section(s) of the report that provide the relevant technical details. The introduction should be aimed at an informed, but otherwise non-expert, reader. A good tip is to assume that all your assessors will read the abstract and introduction, whereas the more detailed technical sections may only be read by your first and second markers - it's therefore really important to get it right.
\end{comment}

This project is follows the consequences of new techniques to optimise
relational queries of joins followed by selections in Haskell. Say you have the
following records in Haskell.

\begin{hscode}\SaveRestoreHook
\column{B}{@{}>{\hspre}l<{\hspost}@{}}%
\column{16}{@{}>{\hspre}c<{\hspost}@{}}%
\column{16E}{@{}l@{}}%
\column{19}{@{}>{\hspre}l<{\hspost}@{}}%
\column{E}{@{}>{\hspre}l<{\hspost}@{}}%
\>[B]{}\mathbf{data}\;\Conid{Students}{}\<[16]%
\>[16]{}\mathrel{=}{}\<[16E]%
\>[19]{}\Conid{S}\;\{\mskip1.5mu \Varid{uid}\mathbin{::}\Conid{Int},\Varid{name}\mathbin{::}\Conid{String},\Varid{age}\mathbin{::}\Conid{Int}\mskip1.5mu\}{}\<[E]%
\\
\>[B]{}\mathbf{data}\;\Conid{Grades}{}\<[16]%
\>[16]{}\mathrel{=}{}\<[16E]%
\>[19]{}\Conid{G}\;\{\mskip1.5mu \Varid{cid}\mathbin{::}\Conid{Int},\Varid{subject}\mathbin{::}\Conid{String},\Varid{grade}\mathbin{::}\Conid{Char}\mskip1.5mu\}{}\<[E]%
\ColumnHook
\end{hscode}\resethooks


\noindent
How might you join two bags of these records together by id? One way would be to
write the following.

\[\lbag\:(s, g)\:|\:s \leftarrow students,\:g \leftarrow grades,\:uid\;s == sid\;g
\:\rbag\]

What are the disadvantages of this approach? Clearly, every possible pair of
elements are considered regardless of whether they have matching ids or not.
This clearly has an asymptotic complexity of $O(nm)$ where $n$ and $m$ are the
cardinality of the $students$ and $grades$ bags respectively.

Databases are vital to modern day society for their ability to structure, sort
and query vast amounts of data from any domain. Of course no one theoretical
model of data has surfaced since its conception, with different approaches
describing how exactly to hold the data. Examples of such data models include
the semi-structured model~\cite{DatabaseSystems} and, more relevant to this
project, the relational model~\cite{RelationalModel} as introduced in
\fref{sec:relationalmodel}. However, all implementations face trade-offs and
development challenges.

This project explores one such solution to the optimisation of queries
consisting of products followed by selections in the relational model. Such
behaviour can be modelled by equijoins, the type of join seen in the example
above, where you wish to keep only two records who share a common attribute.
This project will implement a functioning database system in Haskell according
to the solution outlined in the paper
\relalg{}~\cite{RelationalAlgebraByWayOfAdjunctions}. It then describes original
tools and workflows created to help the benchmarking and testing of the system
in order to evaluate its performance.

Going back to the example, say you could instead index both tables by their
$uid$ and $sid$ respectively. You would not need to consider all possible pair
of elements to find the matching pairs anymore as you would already know which
tuples share the same indexed element. This will allow you to greatly reduce the
order at which you are taking products and, in effect, `making more work for
yourself'. How might this be achieved? The authors of the paper \relalg{},
seemingly enchanted by the elegance of a comprehension notation, have outlined a
category theoretical framework to interpret relational algebra. They found that
many bulk types, such as the bags above, could be made by the adjunction of two
functors and therefore granted them the properties of monads. Moreover, monad
comprehensions have been known about for a substantial amount of
time~\cite{MonadComprehensions} and are generalisation of the expressive list
comprehensions. Therefore, in conjunction with the SQL-like syntax extensions
available to list comprehensions~\cite{ComprehensiveComprehensions} it is hoped
that the defined framework defined in the paper for efficient equijoin can be
expressed in the same concise and powerful manner as the example above.

Continuing from the example above, how might we more efficiently join the
two relations? We could start by indexing both by the key we are looking for and
merging the results.

% TODO: Write code here

After this we simply can perform a Cartesian product on all values matching keys
and reduce the answer back to a bag.

% TODO: Write code here

Note the types of the functions used

% TODO: Write code here

Theoretically, this indexing should help us reduce the asymptotic complexity of
the algorithm. Practically, however, what are the results of such a method? Are
there any limitations? Does it give the gains promised  in real life? This
project is about exploring exactly that idea.

After describing a working implementation of the database system mentioned here,
this report is going to describe the bespoke and original tools created in order
to benchmark the database. We introduce the \relation{JOINBENCH} relation, a
specialised version of past academic benchmarking database schemas made with
testing equijoins in mind. After this, the report will describe the design of
an extensible tool created in order to generate synthetic databases to run the
bechmarks on. This tool can be very easily customised and is made with enforcing
mathematical properties of data in mind so that interesting patterns can emerge
from the \relation{JOINBENCH} relation. Next, an overview of the benchmarking
workflow and experiment designs is given to be open and encourage readers to try
and reproduce results. Finally, a discussion of the results is given and the
efficacy of replacing such a solution is discussed.

To summarise, the main contributions of this project are:
\begin{itemize}
    \item a working implementation of the database system outlined in the paper
        \relalg{},
    \item the definition of the \relation{JOINBENCH} relation specialised from
        old benchmarking databases to provide all the tools needed to test
        equijoin,
    \item a completely new tool to generate any number of general synthetic data
        sets,
    \item an analysis of the results of the benchmarks on the database
        management system described above to evaluate its performance and
        success.
\end{itemize}

\chapter{Background}
\begin{comment}
The background section of the report should set the project into context by relating it to existing published work which you read at the start of the project when your approach and methods were being considered. There are usually many ways of solving a given problem, and you shouldn't just pick one at random. Describe and evaluate as many alternative approaches as possible. The published work may be in the form of research papers found in the academic literature, articles, text books, technical manuals, or even existing software or hardware of which you have had hands-on experience. Your must acknowledge the sources of your inspiration. You are expected to have seen and thought about other people's ideas; your contribution will be putting them into practice in some other context. However, avoid plagiarism: if you take another person's work as your own and do not cite your sources of information/inspiration you are being dishonest. When referring to other pieces of work, cite the sources where they are referred to or used, rather than just listing them at the end. Accidental plagiarism or not knowing how to cite and reference is not a valid reason for plagiarism. Make sure you read and digest the Department's plagiarism document .

In writing the Background chapter you must demonstrate your ability to analyse, synthesise and apply critical judgement. Analysis is shown by explaining how the proposed solution operates in your own words as well as its benefits and consequences. Synthesis is shown through the organisation of your Related Work section and through identifying and generalising common aspects across different solutions. Critical judgement is shown by discussing the limitations of the solutions proposed both in terms of their disadvantages and limits of applicability.

Typically you can look for Background work using different search engines including:

Google Scholar
IEEExplore
ACM Digital Library
Citeseer
Science Direct
Note 1: Often the terms Background, Literature Review, Related Work and State of the Art are used interchangeably.
Note 2: Keyword search is wonderful, but you need the right Keywords.
Note 2: IEEExplore, ACM Digital Library and Science Direct may require you to be on the College network to download the PDF versions of papers. If at home, use VPN.
\end{comment}

% Currently add in only to stop reference error - not ready for submission
\section{The relational model of a database}
We briefly describe the relational model of a database so that we can introduce the key operators we are modelling using category theory. \todo{not modelling, think of better word}
\subsection{Introduction to the relational model}
There are several different data models to choose from when designing a database that specify important aspects of the design such as the structure, operations and constraints on the data \cite{DatabaseSystems}. For this project we concern ourselves with the relational model and its associated algebra.

The relational model represents data through two dimensional tables, called \emph{relations}. Each relation contains several \emph{attributes}, denoting the columns of the table. We call the name of the relation and the set of its attributes a \emph{schema} and they are denoted by the name of the relation followed by the set of attributes in parentheses.\cite{DatabaseSystems} For instance: 
\begin{center}
\begin{verbatim}
  R(a1, a2,..., an)
\end{verbatim}
\end{center}
is a relation \verb|R| with $n$ attributes \verb|a1|, \verb|a2|, \ldots, \verb|an|. A more practical example is the schema
\begin{figure}[!h]
\begin{verbatim}
  People(firstName, surname, age)
\end{verbatim}
\caption[Schema for the People relation]{Example of a schema for the relation People}
\label{fig:peopleSchema}
\end{figure}
This describes a schema for the relation \verb|People| and the attributes \verb|firstName|, \verb|surname| and \verb|age|.
An empty table showing the relation can be found on \fref{tab:peopleRelationHeadings}.
\begin{table}[h]
  \centering
  \begin{tabular}{l|l|l}
    \verb|firstName| & \verb|surname| & \verb|age| \\
    \hline\hline
    & &\\
  \end{tabular}
  \caption{A tabular representation of the People relation.}
  \label{tab:peopleRelationHeadings}
\end{table}
\todo{Add verb to People in the caption}

\subsection{Relational Algebra}
\subsubsection{Projections}
\subsubsection{Selections}
\subsubsection{Products}
\subsubsection{Joins}
\subsection{Evolution of database representation}
\subsubsection{Bags}
\todo{Understand and distinguish between bags being the bulk type}
\paragraph{Characteristics of a database}We expect our database approximation to not be ordered and admit multiplicities and a finite bag of values is one of the simplest constructions that does so. Like a finite set, a bag contains a collection of unordered values. However, unlike a set, bags can contain duplicate elements. \cite{RelationalAlgebraByWayOfAdjunctions} This multiplicity is key for processing non-idempotent aggregations. For instance, if summing up the ages of a database of people, without admitting multiplicity we would only sum each unique age once.
\subparagraph{Generalisation}Furthermore, going forward we generalise to bags of any types instead of the classical ``bags of records''. This also allows us to deal with intermediate tables that contain non-record values.
\todo{Add the mathematical parts about bags}
\todo{When reading about finite maps it says that it's better for databases as non finite maps cannot be aggregated, can non finite bags be aggregates - is this why they are finite?}

\subsubsection{Indexed tables}
We want to move towards an indexed representation of our table in order to equijoin by indexing. \todo{Understand if this is right and equijoin by indexing}. So in this section we introduce the mathematical concepts required to define such an implementation.
\theoremstyle{definition}\newtheorem*{psetdef}{Pointed set}
\theoremstyle{definition}\newtheorem*{ppfuncdef}{Point-preserving function}
\theoremstyle{definition}\newtheorem*{mapdef}{Map}
\theoremstyle{definition}\newtheorem*{finitemapdef}{Finite map}
\begin{psetdef}\label{def:pset}
  A pointed set $\pset{A}{a}$ is a set $A$ with a distinguished element $a \in A$.
\end{psetdef}
We commonly refer to the distinguished element of a set $A$ as $null_A$ or, when not ambiguous to do so, $null$.
\begin{ppfuncdef}\label{def:ppfunc}
  Given two pointed sets $\pset{A}{null_A}$ and $\pset{B}{null_B}$, a total function $f: A \rightarrow B$ is point-preserving if $f(null_A) = null_B$.
\end{ppfuncdef}
\todo{See if there is a way to have better spacing in brackets}

\todo{Add in more mathematical detail for point preserving functions}

We now have the mathematical tools required to define a map. In its finite form a map is widely known in computer science by many other names such as a dictionary, association lists or key-value maps.

Let $\keyset$ be a set and $\valset$ a pointed set. To those already familiar with maps, it may help to think of $\keyset$ as keys and $\valset$ as values.
\begin{mapdef}
  A map of type $\map{\keyset}{\valset}$ is a total function from K to V.
\end{mapdef}
\begin{finitemapdef}
  A finite map of type \finitemap{\keyset}{\valset} is a map where only a finite number of keys are mapped to $null_\valset$ (where $null_\valset$ is the distinguished element of \valset). 
\end{finitemapdef}
The advantage of using a finite map in a database is to allow aggregation.
\todo{Understand why only semi-monoidal}
\todo{Introduced indexed tables} 
\section{Benchmarking databases}
Databases are pervasive to modern society and thus standards have arisen over
the last few decades to ensure that customers are able to pick their preferred
DBMS vendor.

At the top level, database benchmarks are classified into three categories: industry-standard, vendor and
customer-application \cite{PractitionersIntroduction}. 
These classifications are usually motivated by intention of the benchmark
instead of structure of the database management system; there is no shortage of
papers emphasising the importance of domain-specific benchmarks for applications
\cite{PractitionersIntroduction, BenchmarkHandbook} and depending on the
risk/performance tolerance of the application it may be necessary to consider
results from all three categories.

\paragraph{Vendor benchmark} A vendor database benchmark is used by
the database vendor during the production of the database management system.
This benchmark usually serves multiple purposes not just limited to the design
of the system. Of course, it is often used to highlight any performance
bottlenecks driving the design of the system internally but it usually doubles
up and acts as the knowledge base for the marketing of the system. Vendor
benchmarks are usually characterised by a more comprehensive suite of
tests to generate the insights capable of driving the direction of the
product \cite{PractitionersIntroduction}.

\paragraph{Industry-standard benchmark} An industry-standard benchmark is a set
benchmarking suite designed independently to any vendor or solution. It is
designed to allow for a fair comparison between different vendors and has been
shown to increase competition between vendors \cite{Wisconsin2}. As much of
this review will show, many industry-standard benchmarks have been developed
to give results for a wide range of applications of databases for much more
relevant and specialised results; and, although many benchmarks have similar
metrics \cite{SetQueryBenchmark, DebitCredit}, just selecting what metrics to
show consumers is not a trivial task \cite{DebitCredit}. When Gray's paper
\cite{BenchmarkHandbook} was written, it was noted that these benchmarks were
becoming so popular that vendors were also beginning to report their results
with their marketing.

\paragraph{Consumer-application benchmark} This type of benchmark refers to any
benchmarking that a customer would run, typically to choose between different
vendors for their application. This kind of testing can be critical for a
performance-sensitive application \cite{PractitionersIntroduction} and is often
done to test the performance of the database under a specific installation
profile \cite{DoingYourOwnBenchmark} or loads. This is a very specialised
requirements based benchmarking.

The review now turns to different domain-specific benchmarking of databases in
order to comment on and analyse any common structures found when designing
databases. This will help inform major design decisions of this project when
deciding what aspects to include in the benchmarking of the alternative join
syntax.

Despite the innumerable use cases for databases in the modern world, many
applications require and prioritise similar values when choosing a solution.
Furthermore, as most industrial-standard database benchmarks have been designed in
order to allow comparison for specific domain requirements, you will commonly
find database benchmarks that are designed with one of two different
applications in mind: \emph{transaction processing} or \emph{decision support}
\cite{PractitionersIntroduction}; it is worth mentioning that other application
types exist, however, such as \emph{document search} and \emph{direct marketing}
\cite{SetQueryBenchmark}. The application types of transaction processing and
decision support are so pervasive that they are commonly also used to partition
other application types, for instance OLTP (Online transaction processing) and
OLAP (Online analytical processing) are specialised for online processing but
share many comparative similarities as transaction processing and decision
support respectively \cite{OLTP-Oracle}.

This review will provide a brief overview of both
the main database application types, then focus more strongly on decision
support benchmarks as ad-hoc query processing is more heavily tested in these
benchmarks and thus more relevant to the project.

\subsection{Transaction processing}
Transaction processing is characterised by a large number of
update-intensive database services with a particular emphasis on integrity of
its requests \cite{PractitionersIntroduction}.
It is clear that this type of environment
demands an emphasis on throughput and integrity. For instance, a bank would
likely have an online transaction processing system in place and the
ability to deal with a large number of transactions in a short period of
time while maintaining the accuracy of the information is most likely paramount.
There are many well known benchmarks to test these types of systems,
including \emph{DebitCredit}, \emph{TPC-C} and \emph{TPC-E} \cite{DebitCredit,
TPC-OLTP}.
These benchmarks usually measure transactions per second or price per
transactions per second. Efforts have been made to try and describe what is
meant by ``price'' and usually it refers to the five-year capital cost (e.g.
cost of hardware, software and system maintenance over a five year period)
\cite{BenchmarkHandbook, DebitCredit}. The advantage of giving a metric
dependent on price is that price is likely to be a limiting of important factor
to most (if not all) organisations when shopping for database solutions and so
it helps encompass all information about the system into a single comparable
number \cite{SetQueryBenchmark}. It is worth noting that many benchmarks which
include this metric also give more information for the user \cite{TPC-H} or
allow the user to weight the comparative importance of features in this metric
\cite{SetQueryBenchmark} but it is expected that much emphasis will come down to
this defining number \cite{SetQueryBenchmark}.

We note that the above definition of a transaction processing system not only
highlights its intense performance requirements but also notes the importance of
integrity. The bank example highlights intuitively what we mean by integrity, a
shop may not want the transaction to go through zero times if it is reported
that it has, in fact, been processed, just as I would not want the transaction
to go through twice. So what exactly is a transaction? And how can we specify
its integrity in a rigorous way?

A transaction simply is an input to the database that must be considered as one
unit of work \cite{ComputerScienceDictionary} and completed independently to any
concurrent actions occurring \cite{ACID}. The integrity of this behaviour can be described 
by a set of properties typically remembered as the acronym \emph{ACID} and
explained below \cite{ACID, PractitionersIntroduction}.

\paragraph{Atomicity} An ``all or nothing'' property of a transaction. All
individual operations must take place or any partially-completed sequence of
operations must have no lasting effects on the data. The user should be informed
of which outcome has occurred.

\paragraph{Consistency} Given the database was in a consistent/valid state
before the transaction processed, it must remain in a consistent/valid state
after any changes are committed; that is to say the changes cause not integrity
constraints to be violated \cite{IntroToDatabaseSystems}. 
Some texts favour the term ``correctness'' \cite{IntroToDatabaseSystems} but this
discussion is outside of the scope of the project. 
% correctness vs consistency: correctness may be preferred as it is stronger, 
% and says that the database only contains true statements about the world

\paragraph{Isolation} Information related to the execution of one transaction
should be hidden from other transactions running concurrently. This is
equivalent to being able to find an order such that the results of the execution
of all transactions would be the same if they were executed sequentially in that
order.

\paragraph{Durability} Once a transaction has been successfully completed, its
effects must persist and future malfunctions.

\subsection{Decision support}
Decision support is characterised by the need to execute complicated queries on
a database with fewer changes \cite{IntroToDatabaseSystems,
PractitionersIntroduction}. The business value from this sort of application are
vast, especially in automatic report generation to produce insights and direct
company strategy \cite{SetQueryBenchmark, IntroToDatabaseSystems}. Despite the
clear business needs, the measurements usually taken by decision support
benchmarks are not as standard as those by transaction processing benchmarks and
metrics usually vary according to benchmark 
\cite{PractitionersIntroduction}; some common metrics collected are a measure of
elapsed time \cite{Wisconsin, TPC-H, PractitionersIntroduction}, performance of
the underlying system (CPU and I/O utilisation) \cite{PractitionersIntroduction}
and the throughput or price per throughput \cite{TPC-H, SetQueryBenchmark,
PractitionersIntroduction} (as in transaction support systems). 

Although less pervaisive in our every day life to transaction support systems,
there are a large number of benchmarks for decision support systems. To name a
few there are the \emph{Wisconsin benchmark} \cite{Wisconsin}, \emph{Bull Single-User Decision
Support benchmark (SUDS)} \cite{PractitionersIntroduction}, \emph{Set Query
benchmark} \cite{SetQueryBenchmark} and \emph{TPC-H} \cite{TPC-H}. 

\subsection{Database benchmarking case studies}
% Explain why selected these benchmarks

% Wisconsin
% Set query
% DebitCredit

\subsection{Best practices benchmarking}


\begin{comment}
The central part of the report usually consists of three or four chapters detailing the technical work undertaken during the project. The structure of these chapters is highly project dependent. They can reflect the chronological development of the project, e.g. requirements, design, implementation, experimentation, optimisation, evaluation etc. although this is not always the best approach. However you choose to structure this part of the report, you should make it clear how you arrived at your chosen approach in preference to the other alternatives documented in the background. If you have built a new piece of software you should describe and justify your design and details any interesting problems with, or features of, your implementation. Integration and testing are also important to discuss in some cases. You need to discuss the content of these sections thoroughly with your supervisor.
\end{comment}
\section{Synthetic database creation}\label{sec:benchmark:database}
Synthetic databases are critical to running a customisable and repeatable
benchmark. Those who wish to benchmark a database can use custom synthetic
databases in order to specify a large amount of unknowns in the data of the
benchmarks run. Examples of such parameters range from both the size of the
database as a whole to on a tuple level by carefully choosing attribute domains.
Attribute domains can also help give a heuristic control to resource load while
processing queries by determining how difficult it is to computer an answer (for
instance using the equality of strings). Furthermore, when designed correctly synthetic
databases can provide a readable method to easily specify wanted properties of
queries via explicit and easy to calculate selectivity factors in selections and
joins. More information on database benchmark design can be found in
\fref{sec:background:benchmarkbestpractices}.

\paragraph{} In order to have a customisable synthetic benchmark for the
project, we created a collection of python modules to easily configure and
export different databases. We chose to implement this as an extensible set of
classes to represent different cell domains which then can be configured,
specified and passed to objects representing larger structures of the database
and handle generation for each of their responsibilities. This section is a
walk through the ensuing python modules and their internal relationships.

\begin{figure}[t]
\begin{tikzpicture}
    \begin{package}{databasegen}
        \begin{abstractclass}{Cell}{-4,0}
            \operation[0]{+ generate() : str}
        \end{abstractclass}

        \begin{class}{RecordGenerator}{4, 0}
            \attribute{+ cells : List[Cell]}
            \operation{+ generate() : List[str]}
        \end{class}

        \begin{class}[text width=9cm]{TableGenerator}{0, -3}
            \attribute{+ record\_generator : RecordGenerator}
            \operation{+ generate(num\_records : int = 1) : Iterable[List[str]]}
        \end{class}

        \begin{class}[text width=6cm]{CSVTableGenerator}{0, -6}
            \attribute{+ path : str}
            \attribute{+ dialect : str}
            \attribute{+ table\_generator : TableGenerator}
            \operation{+ generate(num\_records : int) : null}
        \end{class}

        \composition{RecordGenerator}{}{}{Cell}
        \composition{TableGenerator}{}{}{RecordGenerator}
        \composition{CSVTableGenerator}{}{}{TableGenerator}
    \end{package}
\end{tikzpicture}
\caption{UML class diagram showing the modules used for synthetic database
creation.}
\label{fig:project:benchmark:creation-class-diagram}
\end{figure}

\paragraph{Overview} This section describes the overall structure of the
solution, a class-diagram to aid the explanation can be found in
\fref{fig:project:benchmark:creation-class-diagram}.
We start with a list of objects extending the \lstinline{Cell}
abstract class. We pass this list to a \lstinline{RecordGenerator} object, which
then can be given to a \lstinline{TableGenerator} object to generate the
database. Finally the database generator can be given to an object, such as the
implemented \lstinline{CSVTableGenerator} to create the
specified output. The advantages of this approach is the ability to easily
configure a handful of properties of the output database by varying just the list of cells.
Cells can be easily manipulated given they conform to the \lstinline{Cell} API
and so their basic capabilities cover a large amount of behaviour and
compositions.

\paragraph{Cell} The \lstinline{Cell} abstract class presents the main API for
customising the domain of attributes in a database table. It defines only one
operation, namely \lstinline[language=Python]{generate() : str} which returns the next value of
the attribute the \lstinline{Cell} class represents. It is worth noting that
once you extend the \lstinline{Cell} abstract class, its method is not static so
an object must be instantiated; this gives every attribute the ability to hold
state to help control the generation of values.

In practice, when extending the cell class for uses within the project, specific
patterns would emerge. Usually, most cells require a
\lstinline{Random} object in their initialiser. The benefit to the composition
with an external \lstinline{Random} object is that the generation of values can
be controlled both in a fine-tuned and systematic way; this allows the user to
specify the distribution used to generate numbers both in a cell scope and
interlink randomness between cells by sharing \lstinline{Random} objects.
Although \lstinline{Cell} is an abstract class, I did not think it would be
appropriate to enforce the use of a \lstinline{Random} object in the initialiser
or give default behaviours. The result of generating a cell can be so varied
that no default behaviour would be appropriate in a general case, potentially
causing future extensions of the \lstinline{Cell} class to behave unexpectedly.
Furthermore, unfortunately python does not support interfaces and favours duck
typing and so abstract classes were a necessary evil to specify a common API and
type the code for static analysis; otherwise an interface would much better suit
the problem domain. As a final addition, abstract classes with behaviour
increases the coupling of the code and so the argument could be made it would
negatively impact the cleanliness of the code.

In order to create the joinbench benchmark a handful of different cells were
designed to represent the domains necessary to create the \relation{JOINBENCH}
table. Great care was taken in the design to eliminate duplication and
generalise the cells as will be presented here.

Most \lstinline{Cell}s used in the \relation{JOINBENCH} table are based on the
\lstinline{RandomModularIntegerCell} class. As the name suggests given a modulo
$n$ in the constructor this cell will generate a random number in the range $\{0,
\,1,\,\ldots,\,n - 1\}$ using a given \lstinline{Random} object. In fact, the
implementation of \lstinline{RandomModularIntegerCell} relies on a more general
\lstinline{RandomIntegerRangeCell} which when given a lower bound $l$ and an
upper bound $u$ generates numbers in the range $\{l,\,l + 1,\,\ldots,\,u - 1\}$.
The \lstinline{RandomIntegerRangeCell} along with the syntactic sugar of the
\lstinline{RandomModularIntegerCell} can represent most attributes seen in the
\relation{JOINBENCH} table; examples are \relationAttribute{onePercent} defined as a
\lstinline{RandomModularIntegerCell} with modulo $100$ or
\relationAttribute{twentyFivePercent} may be implemented with a
\lstinline{RandomIntegerRangeCell} with lower bound $0$ and upper bound $4$
(generating numbers between $0$ and $3$ inclusive for a total of $4$ distinct
values). However, more complex cells that fundamentally are not too different
from the cells described above are needed for a larger set of use cases. How
might we randomly generate entries from the first $10$ values of the geometric
sequence $3^i$ say? Or, more relevant to the \relation{JOINBENCH} relation, the
first $100$ even numbers? To allow a more general number generation from
integers the \lstinline{RandomIntegerRangeCell} and
\lstinline{RandomModularIntegerCell} are based on
\lstinline{RandomModifiedIntegerRangeCell} and
\lstinline{RandomModifiedModularIntegerCell} respectively. How these classes
differ is an extra parameter in the constructor that takes an argument of type
\lstinline[language=Python]{Callable[[int], int]}, the type that represents a function that takes an integer
and returns an integer. Cells of this type generate a random number within the
constraints given by the range or modulo as usual then apply the
\lstinline{modifier} function and return a string representation of the result.
It is clear that the \lstinline{RandomModifiedIntegerRangeCell}
(\lstinline{RandomModifiedModularRangeCell}) and
\lstinline{RandomIntegerRangeCell} (\lstinline{RandomModularIntegerCell}, resp.)
are related by passing the identity function ($\mathrm{id}: x \mapsto x$) as the
\lstinline{modifier} argument. See
\fref{fig:project:benchmark:random-integer-cell-dependencies} for a
comprehensive representation of the dependencies discussed above.

Cells that cover a different generation space or have more complicated
requirements are representable with this structure and design. An
example of a cell that uses and modifies its state, also used in the
\relation{JOINBENCH} relation, is the \lstinline{CounterCell}. The
\lstinline{CounterCell} is used in the \relation{JOINBENCH} relation to provide
a unique integer for each row representing what is commonly known as a primary
key. Its implementation consisting of an integer
attribute initialised in the constructor and returned then incremented after
every call to \lstinline{generate} is a simple advocate to the utility of implementing
\lstinline{Cell} as a class allowing state in contrast to stateless alternatives. Before
\relation{JOINBENCH} was settled on and domains specified other test attributes
for an \database{INVOICE} database were defined and represent a more creative
coverage of possible attribute domains. A \relation{CUSTOMER} table was defined
for the \database{INVOICE} database and it is no surprise a
\relationAttribute{firstName} attribute followed. In response to this attribute
a \lstinline{RandomChoiceCell} was defined which once initialised with a list
of strings (such as names in this example) the \lstinline{generate} method
would randomly pick and return an option from the list. The
\lstinline{RandomChoiceCell} extends the examples presented in the report to
domains other than integers, covering the space of types that can be converted
to strings once the exact domain is enumerated into a list. We also mention the
\lstinline{DateCell} to
demonstrate a more complex operation defined
for use in an \relation{INVOICE} relation within the \database{INVOICE}
database. The \lstinline{DateCell} generated a random valid date contained within the range of
years provided in the constructor returned in long form \verb|YYYYMMDD|.

Before much of the complexity was alleviated by integer based design of the
\relation{JOINBENCH} relation systems had to be design to ensure attributes had
the desired properties. The \relation{JOINBENCH} relation gives us mathematical
approximations, algorithms and laws to ensure properties such as frequency of
values and uniqueness but for the more general space of domains as required by
previous iterations of databases used, such as the \database{INVOICE} database,
alternative methods were required. To fulfil these requirements the idea of cell
composition was introduced. Cell composition is simply a cell that takes another cell
during initialisation to offer a modified behaviour. A prime example of such a
composite cell is the \lstinline{UniqueCell}. As the name suggests the
\lstinline{UniqueCell} ensures that all values generated by the cell are unique.
During initialisation, we might pass an object of type
\lstinline{RandomModularIntegerCell} with modulo $100$ to the
\lstinline{UniqueCell} constructor as follows:
\begin{lstlisting}[language=Python]
nums: Cell = RandomModularIntegerCell(100)
unique_nums: Cell = UniqueCell(nums)
\end{lstlisting}
The result, \lstinline{unique_nums}, is then a cell that generates numbers from
$0$ to $99$ at most once.

\begin{figure}[t]
    \begin{tikzpicture}
        \begin{class}{RandomIntegerRangeCell}{0, -3}
        \end{class}
        \begin{class}[text width=7cm]{RandomModifiedIntegerRangeCell}{0, -0}
        \end{class}
        \begin{class}{RandomModularIntegerCell}{8, -3}
        \end{class}
        \begin{class}[text width=7cm]{RandomModifiedModularIntegerCell}{8, -0}
        \end{class}
        \begin{class}{EvenIntegerCell}{0, 3}
        \end{class}
        \begin{class}{OddIntegerCell}{8, 3}
        \end{class}

        \unidirectionalAssociation{EvenIntegerCell}{}{}{RandomModifiedModularIntegerCell}
        \unidirectionalAssociation{OddIntegerCell}{}{}{RandomModifiedModularIntegerCell}
        \unidirectionalAssociation{RandomIntegerRangeCell}{}{}{RandomModifiedIntegerRangeCell}
        \unidirectionalAssociation{RandomModifiedModularIntegerCell}{}{}{RandomModifiedIntegerRangeCell}
        \unidirectionalAssociation{RandomModularIntegerCell}{}{}{RandomModifiedModularIntegerCell}

    \end{tikzpicture}
    \caption{UML dependency diagram showcasing the web of dependencies affecting
    the random integer group of cells. Note that the APIs have been omitted for
relevance.}
    \label{fig:project:benchmark:random-integer-cell-dependencies}
\end{figure}

\paragraph{RecordGenerator} The \lstinline{RecordGenerator} class is
responsible for defining objects that generate the tuples of a table. As a tuple
is a set of attributes, the \lstinline{RecordGenerator} must be instantiated
with a list of \lstinline{Cell}s. Note that in this implementation the set of
attributes as theoretically outlined in \fref{sec:relationalmodel} gain an
order but this is for the obvious pragmatic use when exporting data. The
\lstinline{RecordGenerator} class defines only one operation,
\lstinline{generate}, but instead of returning a string, it returns a list of
strings in the same order as the list of \lstinline{Cell}s passed to the
constructor.

\paragraph{TableGenerator} The \lstinline{TableGenerator} builds on the
\lstinline{RecordGenerator} in much the same way a table is a set of tuples. It
requires a \lstinline{RecordGenerator} in its constructor to define the domain
and number of attributes and also has a consistently named operation,
\lstinline{generate}. The operation \lstinline{generate} however is subtly
different from the pattern established in \lstinline{RecordGenerator}. Assuming
a table in practice should at least an order of magnitude larger than a single
tuple, it might become costly to generate the entire table at once and hold it
in memory while it is being used. Instead, the \lstinline{TableGenerator}'s
\lstinline{generate} function takes in a parameter describing how many rows to
generate, as one might expect, but returns an \lstinline{Iterable} type where
each iteration returns the list of strings generated by the
\lstinline{RecordGenerator}.

\paragraph{Outputting the data} Although there is no need for a next step,
outputting the data is likely to be a next step. After defining a
\lstinline{TableGenerator} we are left with an function that given the
cardinality of a relation, will product an iterator that generates such a table.
A recommended pattern for outputting data is again to take the
\lstinline{TableGenerator} as an argument during construction and define a
\lstinline{generate} function that specifies a number of records and returns
nothing. This counter intuitive return value is to ensure that the efficiency
gained by the \lstinline{TableGenerator} is not lost as accumulating the table
would cause, where possible it would be beneficial to write the relation one
tuple at a time to whatever output is desired. I did not define an interface as
think that this step in the process is closest to client code and is not needed
by framework outlined therefore it would be needlessly restrictive.
\subparagraph{CSVTableGenerator} An example of such a class used in the project
is the \lstinline{CSVTableGenerator} that follows the exact recommendations
above. It takes a specification of the CSV file to write to, namely a path and
CSV dialect, and generate commands the CSV library to write each record one row
at a time. Of course the operating system may buffer the output but at a high
level this allows for what I believe to fine-grained yet logical writes that the
operating system can schedule.
\subparagraph{Other ideas} There are many other way to potentially output the
data in useful ways. For instance, one may output the database to the standard
output
stream (or perhaps both standard output and a given file). More interestingly one could execute an SQL \lstinline{INSERT}
statement for each tuple, populating a running database.

\section{Benchmark methodology and results}
\subsection{\relation{JOINBENCH} relation and queries}
The \relation{JOINBENCH} relation is a fundamental step in the design of
the benchmark as it singlehandedly determines the possible complexity of the
possible queries. In this regard I think the \relation{JOINBENCH} was mostly
sufficient, it was largely based on previously designed benchmarks with slight
modifications to make it more specialised for the type of interesting
selection and join based queries that would be interesting in this project. One
attribute and family of queries that I found to be missing during my analysis,
however, revolved around the idea of letting the local Cartesian products grow
to be the same size as a normal product on the relations. This could be solved
by adding an attribute \relationAttribute{oneHundredPercent} whose domain only
contained one value (thus containg 100\% of the relation). I believe that this
would more interestingly highlight the disadvantages of the indexed approach as
it is likely to be reduced to a normal product. Admittedly I speculate that the lack of a need
for a filter on $n^2$ elements would still make the indexed equijoin dominant
over the product equijoin; however, I believe it would be much closer to the
comprehension equijoin. Alternatively, this could be done by conducting a
selection on an existing attribute and then joining that attribute to itself,
but I believe this goes against the design decisions expressed in
\fref{sec:background:benchmarkbestpractices}. Taken to the extreme, all
attributes in the relation could be created by various selections on the
\relationAttribute{onePercent} but this not only makes standardising relation
cardinality more difficult, it destroys the readability synthetic data sets are
meant to express. It may be interesting to include attributes in future that
cannot be derived from \relationAttribute{onePercent}, for instance other
multiples such as \relationAttribute{oneThird}, but I am not sure how much more
information this would provide in this case. These suggestion may be taken
forward as future work on designing the \relation{JOINBENCH} relation.

On the topic of readability of attributes, I feel that the naming scheme has
mixed success. The attributes were names after similar attributes in the updated
Wisconsin relation but were kept due to the importance of the cardinality of
partitions to this domain. However, I do find that the attributes names struggle
to easily convey the values in their domains somewhat limiting the expression of
the database. I do feel, however, that if the name gave more information about
the domain it would conversely be more difficult to figure out the partition
sizes. Similarly, all queries are based on uniform sizing and partitions. It
would be a very interesting piece of future work to see how dynamic sized local
joins impact indexing; we could ask questions such as whether a joins with one
partition with a larger cardinality and few other small ones may be faster than
a join with relations with larger tuple counts but many small local partitions.
Having non uniform sizes in the queries would make the benchmark more realistic
to real world scenarios but I also struggle to see a scientific question that
may be asked. Furthermore, I feel the theoretical explanations of the results
obtained in this experiment as described in \fref{sec:benchmark:results}
already gives enough understanding to extrapolate and model into the
non-standard cases.

\subsection{Results}
This subsection will evaluate the results presented in
\fref{sec:benchmark:results} and the processes that obtained them.
At first glance the results of the benchmark seem to be very reasonable and
trends are smooth, expected and clear. A more thorough analysis can be completed
by checking the standard deviations associated with the results. Tables with the
mean and standard deviation for a select number of queries can be found in
\fref{tab:evaluation:std-dev-comparison-onePercent-onePercent},
\fref{tab:evaluation:std-dev-comparison-onePercent-fiftyPercent} and
\fref{tab:evaluation:std-dev-comparison-evenOnePercent-oddOnePercent}. I have
chosen these queries as I believe they most accurately sample the space of
queries and edge cases. I have included the means in this result to give a
comparison of order between the results and spread. It is clear that the
standard deviation in most cases are 3 orders of magnitude smaller than the
mean, potentially indicating that the `quiet' computational environment
(described in \fref{sec:benchmark:experiment}) the
results were gathered in was successful in producing reproducible data. It is
worth noting that the function indexed equijoin continuously seems to have an
order of magnitude higher standard deviation, suggesting a larger spread of
data. This could be due to the fact that the function is more computationally
complex and therefore more sensitive the runtime environment. Most importantly,
due to the nature of such a small spread, it is unlikely that results wrong
relative to each other (no lines overlap) and therefore the conclusions drawn
may be more confidently accepted despite any statistical noise and uncertainty.

\begin{table}[h]
\centering
\begin{tabular}{llll}
\toprule
 & 1000 tuples & 5000 tuples & 10000 tuples \\
\midrule
Product equijoin & 0.0174 ± 7.14e-05 & 0.436 ± 4.54e-05 & 1.72 ± 0.0025 \\
Comprehension equijoin & 0.00848 ± 7.5e-06 & 0.209 ± 7.66e-05 & 0.835 ± 0.000143 \\
Indexed equijoin & 0.00457 ± 5.17e-05 & 0.0242 ± 0.000301 & 0.0855 ± 0.00477 \\
\bottomrule
\end{tabular}

\caption{A table showing the mean time (s) and standard deviation (s) to complete the query `join onePercent and onePercent' for each function.}
\label{tab:evaluation:std-dev-comparison-onePercent-onePercent}
\end{table}

\begin{table}[h]
\centering
\begin{tabular}{llll}
\toprule
 & 1000 tuples & 5000 tuples & 10000 tuples \\
\midrule
Product equijoin & 0.0197 ± 4.09e-05 & 0.49 ± 0.000208 & 1.94 ± 0.00101 \\
Comprehension equijoin & 0.00853 ± 4.38e-06 & 0.217 ± 0.00233 & 0.882 ± 0.00256 \\
Indexed equijoin & 0.0077 ± 0.000159 & 0.0724 ± 0.0017 & 0.358 ± 0.00933 \\
\bottomrule
\end{tabular}

\caption{A table showing the mean time (s) and standard deviation (s) to
complete the query `join onePercent and fiftyPercent' for each function.}
\label{tab:evaluation:std-dev-comparison-onePercent-fiftyPercent}
\end{table}

\begin{table}[h]
\centering
\begin{tabular}{llll}
\toprule
 & 1000 tuples & 5000 tuples & 10000 tuples \\
\midrule
Product equijoin & 0.019 ± 1.71e-05 & 0.477 ± 8.97e-05 & 1.87 ± 0.000181 \\
Comprehension equijoin & 0.00811 ± 2.22e-06 & 0.202 ± 3.31e-05 & 0.791 ± 0.000797 \\
Indexed equijoin & 0.00358 ± 2.23e-05 & 0.00743 ± 6.81e-05 & 0.0153 ± 0.000118 \\
\bottomrule
\end{tabular}

\caption{A table showing the mean time (s) and standard deviation (s) to
complete the query `join evenOnePercent and oddOnePercent' for each function.}
\label{tab:evaluation:std-dev-comparison-evenOnePercent-oddOnePercent}
\end{table}

On the other hand, despite the results having a contained spread, the
conclusions drawn in \fref{sec:benchmark:results} have much space to become more
rigorous. Firstly, a hypothesis test may be conducted to determine whether the
claims are significant in a more prescribed way and given more time I think this
is an important next step. Secondly, much of the section was based on comments
of trends and despite the data clearly being consistent within repeats of the
same query no statistical analysis has been conducted on the trends. Because of
the clear visualisations of the data I am confident that the trends in question
are existent however it would be vital to conduct an analysis to model the
relationships between tuple counts, different queries and functions. This is
clearly where the results fall short and future work would need to be done in
order to add rigour to the already clear patterns the data shows.

Finally, another shortcoming of the results is the potential improvements that
could have been made to statistics considered and ways in which the tools in
question were understood and used. The mean average can easily be swayed by
outliers and, although the spread suggests the data was quite tight and
outliers not too extreme, this could have been a heavy issue. Computer
benchmarking can be prone to outliars for a variety of factors from noise on the
system resources to caching effects and therefore more care should have been
taken to mitigate these risks. However a quick heuristic analysis of the median
and means shows that in this case they are not far apart. Moreover, the
\verb|Criterion| library itself gives a number of more reliable statistics that
I did not use in this project because of time and scope constraints. These may
have led to more reliable results. More time could have been spent on modifying
default options for the library to ensure that sampling rates and other useful
variables were consistent and the benchmarks done in a much fairer way. Luckily,
I believe these did not have a serious impact on the results as the data
gathered seems to be consistent and patterns significantly distinct and clear. I
believe it would also be important to report the usage of other computer
resources, especially CPU utilisation and memory, as these are important
considerations when choosing a query function and standard practice in the academic
database benchmarking community as discussed in
\fref{sec:background:benchmarkbestpractices}.



\chapter{Evaluation}
\chapter{Ethical Issues}
\begin{comment}
You are required to include a short discussion of ethical, legal, societal and professional issues that are relevant to your project (usually 1 to 2 pages long). This should fit in the most appropriate section of your project report, often the Background or Conclusions section.
\end{comment}

\section{Domain ethics}
To my assessment the conclusions of this paper will have a neutral impact on
society. The domain of the research, database management systems, is pervasive in nature and therefore not
applicable to any particular cause.

Databases are a necessity for most modern applications. Furthermore, database
optimisations are especially important for large, scaling products or services.
This means that improving this area of computer science has a very
indiscriminate effect on the type of software products that can make use of it.
Of course, just as any private company could make use of the newer more
efficient techniques, any military hosting a database could too. It is also
worth noting the implementation is in a functional programming language, a
rising paradigm \cite{LanguagePopularity} typically associated with reliability
\cite{RealWorldHaskell} which aligns very well with the real-time and safety critical nature of military applications.

In fact, it could be argued that any findings could benefit society at large in an environmental standpoint. The heart of the project is in database optimisations, which in turn reduce the amount of time required to process queries. One could extrapolate, that without much of an increased processing complexity and resource consumption, this would lead to a reduction in the energy required for these database queries. Given how commonplace database queries are this could have a non-trivial impact.
On the other hand, this research is on a niche subset of database
management systems, only focusing not only on database management systems in
Haskell, a language not very popular for such applications, but on queries
handwritten using list comprehension notation which suggest that any
findings would not be applicable to such scenarios anyway.

\section{Research ethics}
\section{Application}

\chapter{Conclusion}
\begin{comment}
Conclusions and Future Work
The project's conclusions should summarise the key insights that have been gained, be they positive or negative. For example, "The use of overloading in C++ provides a very elegant mechanism for transparent parallelisation of sequential programs", or "The overheads of linear-time n-body algorithms makes them computationally less efficient than O(n log n) algorithms for systems with less than 100000 particles". Avoid tedious personal reflections like "I learned a lot about C++ programming...". It is common to finish the report by listing ways in which the project can be taken further. This might, for example, be a plan for doing the project better if you could do a re-run, turning the project deliverables into a more polished end product, or extending the project into a programme for an MPhil or PhD.
\end{comment}
This chapter is a summary of the findings and original contributions this
project has. The contributions of the project are as follows:
\begin{itemize}
    \item An implementation of the database system described in
        \relalg{}~\cite{RelationalAlgebraByWayOfAdjunctions} in Haskell,
    \item The design of a \relation{JOINBENCH} relation,
    \item A low-level tool for the generation of synthetic databases,
    \item An evaluation of the performance of the aforementioned database
        systems.
\end{itemize}

The beginning of the project was successful with a crude implementation of the
database system described in the paper. The implementation was functional with
unit tests strengthening belief in correctness tests and ease of use during
benchmarks. Slight alterations had to be made to the implementation described by
\relalg{}, a handful of original contributions such as the definition of
equality
of bags were made and much effort was put into the integration of the system to form the
library defined in the code of this project. The library could further be
expanded through the definition of more keys and much real world modification
must be done to the design to make it useful in real world (for instance the
caching of indexed tables). The library defined is usable and functional for the
purposes of benchmarking and a good starting point for further expansion with
much room for many more query optimisations.

The \relation{JOINBENCH} relation is a carefully thought out combination of
benchmarking databases designed by other academics. I
believe that it exceeds at benchmarking equijoin and has a lot more depth than
the scope of this paper. It has attributes designed to investigate the effect of
different orders of products on the performance of equijoins and can be scaled
through select queries on the \relationAttribute{unique} attribute. There is no
doubt as to whether it has achieved its purpose in the scope of this paper.
After spotting trends it may be extended with a few more attributes or
supporting queries to investigate edge cases more thoroughly.

The synthetic database generator gives users a low-level tool to generate very
customisable databases. The framework is very flexible yet prescriptive of
structure. For the benchmarking of this project, a wide array of different cell
types were defined and used to create the \relation{JOINBENCH} relation and
previous iterations of benchmarking databases. It has a capacity for both
qualitative and quantitative data, however I have found it extremely useful on the
quantitative side. Patterns and trends between cell values can be moderated
through cell composition or passing shared objects into constructors and is very
powerful for customisation of the result. In future this framework could be
extended to allow for a more user-friendly interaction, though care must be
taken as to not lose the flexibility that comes from the low-levelled nature of
the design. Furthermore, work may be made to make the framework more context
aware so that it may be used to generate more useful tools such as SQL commands
to create the database and Haskell parsers to read its output; similarly, care
must be taken not to lose the simplicity of the current solution and force users
to give irrelevant information to their use case. This framework is a large
original contribution of this project.

Finally, although there is much room for more rigour in the statistical analysis
of results, benchmarking with the implementation above and this methodology has
produced extremely strong trends. It is clear that, as expected, indexed
equijoin needs to have a substantial number of tuples in order to become
effective over methods with a worse asymptotic complexity but reduced overhead.
Furthermore, comprehension equijoin does much better than expected and although
scales similar to product equijoin, does so with a significantly lower slope.
The trends behind indexed equijoin are more complicated, however, and rely
heavily on the number of and order of local products needed by the query. This
means that the performance increase of indexed equijoin heavily relies on the
query in question, though most real world situations would benefit from its
implementation as they are unlikely to fall into the same category as edge
cases. On the other hand, there are many query optimisations (such as a
preceding join) that are not considered in this project and future work may be
done to investigate how much of an effect these have on the efficiency of the
solutions. Furthermore, more work can be done to investigate the effect of
non-uniformly sized relations on such queries.

In conclusion, all the tools described in this paper were paramount to the
evaluation of the algorithms presented in the paper \relalg{}. Together they
produced results with relatively low standard deviations within their own
samples and clearly strong trends that are unlikely to be variable. Although
more work can be done to investigate the effect different contexts have on the
results I am confident that the findings of this paper will inform a strong
theoretical background and help inform any conjectures.


% Back matter
\bibliographystyle{unsrtnat}
% DoC uses the Vancouver Referencing Format.
\bibliography{report/bibs/combined}

\end{document}